\section{Quadrilateral Mesh}
A conforming mesh of quadrilaterals will be used in both
finite element and finite volume computations.
Discretization on a quadrilateral mesh require fewer elements than 
discretization on a triangular mesh. Each quadrilateral element is
assumed to be strictly convex, not degenerate into a triangle, line
or point. We use a logically rectangular mesh for the ease of
implementation, which can be counted with 
cartesian index $\{i,j\}$. We show a square element, which is a special
case of quadrilateral mesh, and a general quadrilateral mesh in 
Figure~\ref{fig:meshexample}.
\begin{figure}[ht]
\centering
\input chapters/tikz/quadmesh.tex
\caption{An exhibition of a square and a general quadrilateral mesh.}
\label{fig:meshexample}
\end{figure}
There exists a bijective and smooth map $\mathbf{F}_{\hat{Q}}: \hat{Q} \rightarrow Q$.
